\section{Données du poste guichet (contrôleur PFC200)}
\label{secDonneesGuichet}

Le contrôleur PFC200, scruté sur le réseau \texttt{NOC\_ETHIP\_CONVOYEUR},
se comporte comme un simple îlot d'entrées/sorties pour la commande du
poste guichet. Les données qu'il produit et consomme sont de type
\texttt{BOOL}, obtenues par un service d'IO-Scanning : elles sont non
forçables et leur dynamique (front) n'est pas accessible.

\subsection{Données produites (T->O) -- capteurs du guichet}

\begin{center}
\begin{tabular}{|l|l|p{7cm}|}
\hline
\textbf{Capteur} & \textbf{Référence} & \textbf{Information portée} \\
\hline
aid & \texttt{nixAid} & Capteur validant l'information du capteur de
détection de gobelet \\
\hline
dgt & \texttt{nixDgt} & Capteur de détection de gobelet (information valide
uniquement si \texttt{aid} est actif) \\
\hline
dpp & \texttt{nixDpp} & Capteur de détection d'une palette en poste \\
\hline
dpo & \texttt{nixDpo} & Capteur de détection que le poste est occupé \\
\hline
\end{tabular}
\end{center}

\begin{UPSTIremarque}[Codage de l'état de la palette]
Le couple (aid, dgt) code l'état de la palette détectée : $(1,0)$ signifie
que la palette est vide, $(1,1)$ signifie qu'elle porte un gobelet.
\end{UPSTIremarque}

Ces quatre informations occupent les bits 0 à 3 de l'octet~0 de la zone de
données produites (\texttt{Guichet\_IN}). Les octets 1 à 3 sont réservés
pour une utilisation future.

\begin{UPSTIremarque}[Durée d'immobilisation des palettes pleines]
Les octets 4 à 7 de cette même zone portent une donnée de type
\texttt{TIME} codée en Little Endian, et les octets 8 à 11 la même
information codée en Big Endian : il s'agit de la durée d'immobilisation
des palettes pleines dans le guichet (valeur d'exemple : \texttt{T\#4s}).
\end{UPSTIremarque}

\subsection{Données consommées (O->T) -- actionneurs du guichet}

\begin{center}
\begin{tabular}{|l|l|p{7cm}|}
\hline
\textbf{Actionneur} & \textbf{Référence} & \textbf{Nature} \\
\hline
BE & \texttt{nqxBE} & Abaisser le bloqueur d'entrée (distributeur
monostable) \\
\hline
BP & \texttt{nqxBP} & Abaisser le bloqueur du poste (distributeur
monostable) \\
\hline
CR & \texttt{nqxCR} & Colonne lumineuse RVB -- composante rouge \\
\hline
CV & \texttt{nqxCV} & Colonne lumineuse RVB -- composante verte \\
\hline
CB & \texttt{nqxCB} & Colonne lumineuse RVB -- composante bleue \\
\hline
\end{tabular}
\end{center}

Ces cinq commandes occupent les bits 0 à 4 de l'octet~0 de la zone de
données consommées (\texttt{Guichet\_OUT}). Les octets 1 à 3 sont réservés
pour une utilisation future.

\begin{UPSTIremarque}[Adresses mémoire non données]
Les adresses \texttt{\%MW} de base de ces deux zones ne sont pas données
ici : leur détermination, à partir de la cartographie mémoire des
coupleurs \texttt{BMX~NOC~0401.2} (cf. document
\texttt{02c\_configuration\_noc}), fait partie du travail demandé aux
étudiants.
\end{UPSTIremarque}

\subsection{Volumétrie globale}

Le contrôleur PFC200 produit 12 octets et consomme 4 octets, seul le
premier octet de chaque zone étant significatif pour la commande du poste
guichet.
